\section{Ordered physical histories and repeated access}
\label{sec:new-histories}
These extensions follow the route and coherent-record join. An ordered
word is assigned a unitary only after its physical event/route/control
incidences have been specified. For such unitaries on the one full marker
space define
\[\begin{gathered}
G_{wv}=\operatorname{Tr}(\rho U_v^*U_w),\\
 (\rho_{route})_{wv}=a_w a_v^*G_{wv}.
\end{gathered}\]
The matrix G is positive semidefinite with diagonal one: purification
expresses it as the overlap matrix of normalized conditional marker
vectors (with the indicated index convention). Thus the route density
operator is normalized and positive for normalized route amplitudes.
Every multi-history interference measurement follows from that same
operator and the stated analyser. On a declared accessible domain the
conditional states are \(\mathcal A(U_w\rho U_w^*)\); equality of these
states means identical probabilities for its marker-only POVMs.
It does not imply equality of the coherent cross blocks, nor does a
computational history label alone supply an extra observable.
Equal signed action and final address therefore leave a concrete
question about the ordered physical propagators rather than an
unassigned rule creating additional states. In the moving-source model,
\([H_0,P_f]\) need not vanish as a formal core commutator, so ordered
velocity controls need not be determined by their integrals. The actual
form propagators, not that potentially unbounded commutator, define
the history comparison.

\subsection{A finite repeated-read theorem with correlated noise}
Let Q0,Q1 be the orthogonal populations of the retained memory. An ideal
copy-read isometry has
\(W|j\rangle_R|0\rangle_A=|j\rangle_R|j\rangle_A\), extended
unitarily on its complete finite space. A finite positive controlled
flip realizes this operation with area \(\pi\), peak \(\pi/\tau\)
and absolute switch-work at most \(2\pi/\tau\); the memory/bath drift
must be charged as in the finite pulse construction.
Diagonal interactions conserve Q populations even with a correlated
bath. For transverse disturbance \(V_\perp\), retain each interval's
trace-distance exposure \(d_k\leq\int\|V_\perp(t)\|dt\).
Suppose actual read k differs from the ideal nondestructive read by
\(\delta_k\) on the complete reachable joint state, including all
earlier read registers and the bath. Successive hybrid replacement and
trace-distance contraction give
\[
 \begin{gathered}\Pr\left\{\begin{gathered}\text{any wrong stored/read label}\\\text{through read }N\end{gathered}\right\}\\
 \leq\min\{1,e_{write}+\sum_{k=1}^N(d_k+\delta_k)\}.\end{gathered}
\]
No statistical independence is used. For worst-label uniform bounds
\(d_k+\delta_k\leq b\), an allowed total error \(e_{max}\) permits
\(N b\leq e_{max}-e_{write}\); this is a finite guaranteed-read budget,
not a derived thermal decay rate. Actual sequential joint-state bounds
are essential: separate one-time marginal read errors do not establish
this premise. Classical repetition improves reliability exponentially
only under its additional conditional-independence assumptions.

Every read needs a recipient. An occupied register cannot be reused as
a blank without an explicit reversible transfer or disposal/reset channel.
Exact discrimination of N complete histories requires N orthogonal
recipient states; the integrated approximate label-capacity statement
gives its finite-error version. Resetting the local cell does not erase
its globally retained source/bath/history recipients. A thermal heat
claim additionally needs a specified reservoir, initial correlations and
cyclic protocol. No universal heat-per-Write number is assigned here.

